\chapter{The Operator Product Expansion}
\label{ch:volume-iii-operator-product-expansion}

\section{Local Expansion in Radial Quantization}

Let \(\mathcal O_i(x)\) and \(\mathcal O_j(0)\) be local operators inserted
inside a sphere of radius \(R\), and let all other insertions lie outside that
sphere.  Radial quantization turns the product
\[
  \mathcal O_i(x)\mathcal O_j(0)\ket{\vac},
  \qquad |x|<R,
\]
into a state on \(S_R^{D-1}\).  Since the Hilbert space on the sphere is
spanned by local-operator states, this state can be expanded in a basis of
primary multiplets:
\begin{equation}
  \mathcal O_i(x)\mathcal O_j(0)
  =
  \sum_{\mathcal O}
  C_{ij\mathcal O}(x,\partial)\,\mathcal O(0).
  \label{eq:ope-general-local}
\end{equation}
The sum runs over primary operators \(\mathcal O\) and the differential
operator \(C_{ij\mathcal O}(x,\partial)\) generates the descendants of
\(\mathcal O\).  Equation \eqref{eq:ope-general-local} is understood inside
correlation functions with all other insertions farther from the origin than
\(|x|\), or equivalently as a Hilbert-space expansion on the separating
sphere.

Equivalently, write \(x=\ee^\tau n_x\) under the Weyl map
\(\R^D\setminus\{0\}\simeq\R_\tau\times S^{D-1}\).  For scalar primaries,
\(\mathcal O_i(x)=|x|^{-\Delta_i}\widetilde{\mathcal O}_i(\tau,n_x)\),
so the product creates the cylinder state
\[
  \widetilde{\mathcal O}_i(\tau,n_x)\ket{\mathcal O_j}
  \in\mathcal H_{S^{D-1}}.
\]
The OPE is the statement that this state is represented, in every exterior
correlation function, by a convergent limit of local operators at the origin,
organized into primary states and their descendants.  Thus the expansion is
not a formal Taylor series in elementary fields; it is the spectral expansion
of a local state in the Hilbert space of the conformal theory.

The convergence parameter is radial separation.  If the nearest outside
insertion is at radial distance \(R\), the expansion is organized by powers of
\(|x|/R\), with larger scaling dimensions contributing at higher radial
energy.  Thus the OPE is a structural expansion controlled by the local
operator spectrum and the conformal representation theory of descendants.

\begin{theorem}[Radial OPE convergence]
  Assume Hypothesis~\ref{hyp:radial-reconstruction-data}.  Let
  \(\mathcal O_i(x)\mathcal O_j(0)\vac\) be inserted inside a sphere of
  radius \(r\), and let all remaining insertions lie outside radius \(R>r\).
  Then the expansion of this state in eigenspaces of the dilatation operator
  converges in Hilbert norm for \(r<R\), and the resulting OPE converges inside
  the corresponding correlation functions.  Truncating to states of scaling
  dimension at most \(\Delta_*\) gives a remainder controlled by powers of
  \(r/R\) determined by the omitted dimensions.
\end{theorem}

The proof is the spectral theorem applied to radial quantization.  Radial time
is \(\tau=\log |x|\), and the dilatation operator is the Hamiltonian
generating translations in \(\tau\).  A separating sphere at radius \(R\)
inserts the damping factor \(\exp[-(\log(R/r))D_{\rm cyl}]\), where
\(D_{\rm cyl}\) is the cylinder Hamiltonian.  Since the spectrum is discrete
with finite multiplicities below any fixed dimension, the spectral projection
onto \(\Delta\le\Delta_*\) gives a finite partial OPE and the complement is
suppressed by the corresponding power of \(r/R\) relative to the norm of the
external state.  This is the mathematical reason the OPE is an operator
algebraic expansion in unitary CFTs satisfying the stated radial hypotheses.

\begin{figure}[H]
  \centering
  \resizebox{0.94\textwidth}{!}{%
  \begin{tikzpicture}[x=1cm,y=1cm,every node/.style={font=\scriptsize}]
    \tikzset{
      insertion/.style={circle,fill=violet!80!black,inner sep=1.7pt},
      arrow/.style={-{Latex[length=2mm]},line width=0.8pt},
      surface/.style={green!45!black,line width=0.8pt}
    }
    \begin{scope}
      \node[font=\small] at (0,2.0) {two nearby insertions};
      \draw[thick] (0,0) circle (1.45);
      \draw[surface] (0,0) circle (0.92);
      \node[insertion] at (-0.28,0.10) {};
      \node[insertion] at (0.35,0.35) {};
      \node[violet!80!black] at (-0.58,-0.20) {\(\mathcal O_j(0)\)};
      \node[violet!80!black] at (0.82,0.55) {\(\mathcal O_i(x)\)};
      \node[green!45!black] at (1.12,-0.70) {\(S_R^{D-1}\)};
      \node[blue!65!black] at (0,-1.92) {outside insertions are outside \(S_R^{D-1}\)};
    \end{scope}

    \draw[arrow] (2.00,0.0)--(3.25,0.0);
    \node[align=center] at (2.62,0.52) {expand\\on sphere};

    \begin{scope}[shift={(5.55,0)}]
      \node[font=\small] at (0,2.0) {primary families};
      \draw[thick] (0,0) circle (1.45);
      \draw[surface] (0,0) circle (0.92);
      \foreach \y/\lab in {0.66/{\mathcal O},0.22/{\partial\mathcal O},-0.22/{\partial^2\mathcal O},-0.66/{\cdots}} {
        \node[insertion] at (0,\y) {};
        \node[violet!80!black,right] at (0.18,\y) {\(\lab\)};
      }
      \node[orange!85!black] at (0,-1.85)
        {\(\displaystyle \sum_{\mathcal O} C_{ij\mathcal O}(x,\partial)\mathcal O(0)\)};
    \end{scope}
  \end{tikzpicture}
  }
  \caption{The OPE is the expansion of the state created by nearby insertions
  in the basis of primary conformal multiplets on a separating sphere.}
  \label{fig:ope-radial-expansion}
\end{figure}

\section{Scalar-Scalar OPE}

Let \(\phi_1\) and \(\phi_2\) be scalar primaries of dimensions
\(\Delta_1\) and \(\Delta_2\).  Let
\(\mathcal O_{\mu_1\cdots\mu_\ell}\) be a symmetric traceless primary of spin
\(\ell\) and dimension \(\Delta\).  Its leading contribution to the
\(\phi_1\times\phi_2\) OPE has the form
\begin{equation}
  \phi_1(x)\phi_2(0)
  \supset
  \lambda_{12\mathcal O}
  \frac{x^{\mu_1}\cdots x^{\mu_\ell}}
       {|x|^{\Delta_1+\Delta_2-\Delta+\ell}}
  \mathcal O_{\mu_1\cdots\mu_\ell}(0)
  +\text{descendants}.
  \label{eq:scalar-scalar-leading-ope}
\end{equation}
The tensor \(x^{\mu_1}\cdots x^{\mu_\ell}\) selects the symmetric traceless
part of the spin-\(\ell\) representation.  Antisymmetric components do not
appear in the product of powers of a single vector, and trace parts are
absorbed into descendants or lower-spin structures.

For an exchanged scalar primary \(\mathcal O\), the first descendant term is
fixed by conformal symmetry:
\begin{equation}
  \phi_1(x)\phi_2(0)
  \supset
  \lambda_{12\mathcal O}
  |x|^{\Delta-\Delta_1-\Delta_2}
  \left(
    1+
    \frac{\Delta+\Delta_1-\Delta_2}{2\Delta}
    x^\mu\partial_\mu+\cdots
  \right)\mathcal O(0).
  \label{eq:scalar-ope-first-descendant}
\end{equation}
The coefficient of the derivative is determined by requiring the OPE to
reproduce the conformally fixed three-point function
\(\langle\phi_1\phi_2\mathcal O\rangle\).  To see the coefficient, place the
third operator at \(z\) with \(|x|\ll |z|\) and normalize
\(\langle\mathcal O(y)\mathcal O(z)\rangle=|y-z|^{-2\Delta}\).  The
three-point function gives
\[
  \langle\phi_1(x)\phi_2(0)\mathcal O(z)\rangle
  =
  \lambda_{12\mathcal O}
  |x|^{\Delta-\Delta_1-\Delta_2}|z|^{-2\Delta}
  \left[
    1+
    (\Delta+\Delta_1-\Delta_2)
    \frac{x\cdot z}{z^2}
    +O(x^2)
  \right].
\]
The OPE ansatz
\[
  \lambda_{12\mathcal O}
  |x|^{\Delta-\Delta_1-\Delta_2}
  \bigl(1+a\,x^\mu\partial_\mu+\cdots\bigr)\mathcal O(0)
\]
gives instead
\[
  \lambda_{12\mathcal O}
  |x|^{\Delta-\Delta_1-\Delta_2}|z|^{-2\Delta}
  \left[
    1+2a\Delta\frac{x\cdot z}{z^2}+O(x^2)
  \right],
\]
because
\(\partial_{y^\mu}|y-z|^{-2\Delta}|_{y=0}
=2\Delta z_\mu |z|^{-2\Delta-2}\).  Matching the two expansions gives
\[
  a=\frac{\Delta+\Delta_1-\Delta_2}{2\Delta}.
\]
Higher descendants are fixed in the same way by matching higher orders, or
equivalently by the conformal algebra.

For identical scalar primaries this reduces to the especially important
normalization
\begin{equation}
  \phi(x)\phi(0)
  \supset
  \lambda_{\phi\phi\mathcal O}
  |x|^{\Delta_{\mathcal O}-2\Delta_\phi}
  \left(1+\frac12 x^\mu\partial_\mu+\cdots\right)\mathcal O(0)
  \label{eq:identical-scalar-scalar-ope-first-descendant}
\end{equation}
for scalar exchange.  The coefficient \(1/2\) is not an extra convention; it
is the specialization of the conformal Ward constraints to two equal external
dimensions.

\section{OPE Data}

Choose a basis of primaries with diagonal two-point functions.  The local CFT
data are then
\begin{equation}
  \left\{
    \Delta_{\mathcal O},\;
    \rho_{\mathcal O},\;
    \lambda_{ij\mathcal O}^{(a)}
  \right\}.
  \label{eq:cft-local-ope-data}
\end{equation}
Here \(\rho_{\mathcal O}\) is the spin representation, and the label \(a\)
enumerates independent three-point tensor structures when more than one
structure exists.  Once this data is specified, all descendant coefficients in
the OPE are representation-theoretic consequences of the conformal algebra.

In representation language, if \(V_i\) is the conformal multiplet generated
by a primary \(\mathcal O_i\), the OPE map is a conformally covariant
intertwining map from the local product of two multiplets to the Hilbert space
generated by local operators.  On the decomposition into irreducible conformal
multiplets this has the form
\[
  V_i\otimes V_j
  \longrightarrow
  \bigoplus_{\mathcal O} m_{ij}^{\mathcal O}V_{\mathcal O}
  \subset \mathcal H_{S^{D-1}} .
\]
The multiplicity \(m_{ij}^{\mathcal O}\) equals the number of independent
three-point structures connecting \(i\), \(j\), and \(\mathcal O\).  When
\(\mathcal O_i\) and \(\mathcal O_j\) are both scalars, rotational covariance
permits only symmetric traceless tensor primaries in the leading term, and
the multiplicity for each such representation is one.  For spinning external
operators, several independent tensor structures may couple the same three
representations, and the OPE coefficients must then carry the corresponding
structure label.

Repeated OPEs determine every separated Euclidean correlation function from
the primary spectrum, the two-point pairing, and the three-point
coefficients, provided the radial convergence assumptions stated above hold.
The data are therefore local in a strong sense: higher-point functions do not
bring in independent local constants.  Their nontrivial content is the
compatibility of the same OPE data in different channels.

\section{Conformal Blocks as Multiplet Contributions}

For a scalar four-point function, the OPE separates the contribution of each
primary conformal multiplet.  For the general scalar four-point function
\eqref{eq:scalar-four-point-general}, the \(12\to34\) OPE gives
\begin{equation}
  \mathcal G(u,v)
  =
  \sum_{\mathcal O,a,b}
  \lambda_{12\mathcal O}^{(a)}
  N_{\mathcal O}^{ab}
  \lambda_{34\mathcal O}^{(b)}
  G_{\Delta,\ell}^{(a,b;12|34)}(u,v).
  \label{eq:general-scalar-conformal-block-expansion}
\end{equation}
Here \(\mathcal O\) runs over primary multiplets of dimension \(\Delta\) and
spin \(\ell\), \(a,b\) label possible tensor structures in the
\(\mathcal O_1\mathcal O_2\mathcal O\) and
\(\mathcal O_3\mathcal O_4\mathcal O\) three-point functions, and
\(N_{\mathcal O}^{ab}\) is the inverse two-point pairing on the corresponding
degeneracy space.  In an orthonormal basis and for scalar external operators
with a single structure, this reduces to
\[
  \mathcal G(u,v)
  =
  \sum_{\mathcal O}
  \lambda_{12\mathcal O}\lambda_{34\mathcal O}\,
  G_{\Delta,\ell}^{(12|34)}(u,v).
\]
The function \(G_{\Delta,\ell}^{(12|34)}\) is the conformal block: it is the
sum of the primary \(\mathcal O\) and all of its descendants, with descendant
coefficients fixed by the conformal algebra.

\begin{figure}[H]
  \centering
  \resizebox{0.92\textwidth}{!}{%
  \begin{tikzpicture}[x=1cm,y=1cm,every node/.style={font=\scriptsize}]
    \tikzset{
      insertion/.style={circle,fill=violet!80!black,inner sep=1.7pt},
      line/.style={line width=0.8pt},
      arrow/.style={-{Latex[length=2mm]},line width=0.8pt}
    }
    \begin{scope}
      \node[font=\small] at (0,2.05) {\(12\to34\) channel};
      \node[insertion] at (-2.2,1.0) {};
      \node[insertion] at (-2.2,-1.0) {};
      \node[insertion] at (2.2,1.0) {};
      \node[insertion] at (2.2,-1.0) {};
      \node[violet!80!black,left] at (-2.2,1.0) {\(\mathcal O_1\)};
      \node[violet!80!black,left] at (-2.2,-1.0) {\(\mathcal O_2\)};
      \node[violet!80!black,right] at (2.2,1.0) {\(\mathcal O_4\)};
      \node[violet!80!black,right] at (2.2,-1.0) {\(\mathcal O_3\)};
      \draw[line] (-2.2,1.0)--(-0.75,0);
      \draw[line] (-2.2,-1.0)--(-0.75,0);
      \draw[line,blue!65!black] (-0.75,0)--(0.75,0);
      \draw[line] (0.75,0)--(2.2,1.0);
      \draw[line] (0.75,0)--(2.2,-1.0);
      \node[blue!65!black,above] at (0,0) {\(\mathcal O\)};
      \node[orange!85!black] at (-1.12,-0.34) {\(\lambda_{12\mathcal O}\)};
      \node[orange!85!black] at (1.10,-0.34) {\(\lambda_{34\mathcal O}\)};
    \end{scope}
    \draw[arrow] (3.35,0)--(4.65,0);
    \node at (4.00,0.48) {same correlator};
    \begin{scope}[shift={(7.05,0)}]
      \node[font=\small] at (0,2.05) {\(14\to23\) channel};
      \node[insertion] at (-2.2,1.0) {};
      \node[insertion] at (-2.2,-1.0) {};
      \node[insertion] at (2.2,1.0) {};
      \node[insertion] at (2.2,-1.0) {};
      \node[violet!80!black,left] at (-2.2,1.0) {\(\mathcal O_1\)};
      \node[violet!80!black,left] at (-2.2,-1.0) {\(\mathcal O_2\)};
      \node[violet!80!black,right] at (2.2,1.0) {\(\mathcal O_4\)};
      \node[violet!80!black,right] at (2.2,-1.0) {\(\mathcal O_3\)};
      \draw[line] (-2.2,1.0)--(-0.20,0.72);
      \draw[line] (2.2,1.0)--(-0.20,0.72);
      \draw[line,blue!65!black] (-0.20,0.72)--(0.20,-0.72);
      \draw[line] (-2.2,-1.0)--(0.20,-0.72);
      \draw[line] (2.2,-1.0)--(0.20,-0.72);
      \node[blue!65!black,right] at (0.18,0) {\(\mathcal O'\)};
      \node[orange!85!black] at (-0.65,1.10) {\(\lambda_{14\mathcal O'}\)};
      \node[orange!85!black] at (-0.55,-1.16) {\(\lambda_{23\mathcal O'}\)};
    \end{scope}
  \end{tikzpicture}
  }
  \caption{A four-point function may be expanded in different OPE trees.  The
  equality of these decompositions is the four-point form of OPE
  associativity.}
  \label{fig:ope-crossing-trees}
\end{figure}

The block is characterized representation-theoretically by the quadratic
Casimir.  Let \(\mathcal C_{12}\) be the differential operator representing
the quadratic conformal Casimir acting on points \(x_1,x_2\).  Then
\[
  \mathcal C_{12}\,
  G_{\Delta,\ell}^{(12|34)}(u,v)
  =
  C_{\Delta,\ell}\,
  G_{\Delta,\ell}^{(12|34)}(u,v),
\]
with eigenvalue
\[
  C_{\Delta,\ell}
  =
  \Delta(\Delta-D)+\ell(\ell+D-2).
\]
The boundary condition is fixed by the leading OPE term as \(x_1\to x_2\).
Thus the spectrum and OPE coefficients are the dynamical data, while the
blocks are kinematic functions determined by conformal representation theory.

\subsection{Radial Expansion of a Block}

For identical Hermitian scalar primaries \(\phi\), there is a particularly
transparent block expansion in the radial variable \(\rho\).  Put the four
points in a plane at
\[
  x_1=\rho,\qquad x_2=-\rho,\qquad x_3=-1,\qquad x_4=1,
  \qquad
  \rho=r\ee^{\ii\alpha},\quad 0\le r<1.
\]
With the cross-ratio convention of
\eqref{eq:cross-ratios-u-v}, this frame gives
\begin{equation}
  z=\frac{4\rho}{(1+\rho)^2},
  \qquad
  \rho=\frac{z}{(1+\sqrt{1-z})^2}.
  \label{eq:rho-cross-ratio-map}
\end{equation}
The Weyl map to the cylinder places the two inner insertions at radial time
\(\tau=\log r\) and angular positions separated by \(\pi\), while the two
outer insertions lie at \(\tau=0\).

\begin{figure}[H]
  \centering
  \resizebox{0.88\textwidth}{!}{%
  \begin{tikzpicture}[x=1cm,y=1cm,every node/.style={font=\scriptsize}]
    \tikzset{
      insertion/.style={circle,fill=violet!80!black,inner sep=1.7pt},
      arrow/.style={-{Latex[length=2mm]},line width=0.8pt},
      guide/.style={blue!65!black,dashed,line width=0.75pt}
    }
    \begin{scope}
      \node[font=\small] at (0,2.05) {plane};
      \draw[arrow] (-2.0,0)--(2.0,0);
      \draw[arrow] (0,-1.45)--(0,1.55);
      \node[insertion] at (-1.35,0) {};
      \node[violet!80!black,below] at (-1.35,0) {\(-1\)};
      \node[insertion] at (1.35,0) {};
      \node[violet!80!black,below] at (1.35,0) {\(1\)};
      \draw[orange!85!black,line width=0.8pt] (0,0)--(0.70,0.55);
      \draw[orange!85!black,line width=0.8pt] (0,0)--(-0.70,-0.55);
      \node[insertion] at (0.70,0.55) {};
      \node[insertion] at (-0.70,-0.55) {};
      \node[violet!80!black,above right] at (0.70,0.55) {\(\rho\)};
      \node[violet!80!black,below left] at (-0.70,-0.55) {\(-\rho\)};
      \draw[guide] (0.58,0) arc[start angle=0,end angle=38,radius=0.58];
      \node[blue!65!black] at (0.78,0.20) {\(\alpha\)};
    \end{scope}

    \draw[arrow] (2.75,0)--(4.0,0);
    \node[align=center] at (3.38,0.55) {radial\\quantization};

    \begin{scope}[shift={(6.35,0)}]
      \node[font=\small] at (0,2.05) {cylinder};
      \draw[thick] (-1.45,-1.35)--(-1.45,1.35);
      \draw[thick] (1.45,-1.35)--(1.45,1.35);
      \draw[thick] (0,1.35) ellipse (1.45 and 0.35);
      \draw[thick] (0,-1.35) ellipse (1.45 and 0.35);
      \draw[guide] (0,0.70) ellipse (1.45 and 0.35);
      \draw[guide] (0,-0.25) ellipse (1.45 and 0.35);
      \node[insertion] at (1.45,0.70) {};
      \node[insertion] at (-1.45,0.70) {};
      \node[violet!80!black,right] at (1.45,0.70) {\((0,0)\)};
      \node[violet!80!black,left] at (-1.45,0.70) {\((0,\pi)\)};
      \node[insertion] at (0.65,-0.05) {};
      \node[insertion] at (-0.65,-0.45) {};
      \node[violet!80!black,right] at (0.65,-0.05)
        {\((\log r,\alpha)\)};
      \node[violet!80!black,below left] at (-0.65,-0.45)
        {\((\log r,\alpha+\pi)\)};
      \draw[green!45!black,<->,line width=0.8pt] (-1.12,-0.25)--(-1.12,0.70);
      \node[green!45!black,right] at (-1.12,0.20) {\(-\log r\)};
    \end{scope}
  \end{tikzpicture}
  }
  \caption{The radial variable \(\rho=r\ee^{\ii\alpha}\) puts the
  \(12\to34\) OPE channel into a cylinder frame.  Inserting a complete set of
  cylinder energy eigenstates gives a convergent expansion in powers of \(r\).}
  \label{fig:rho-cylinder-frame}
\end{figure}

Assume \(\phi\) is normalized by
\(\langle\!\langle\phi|\phi\rangle\!\rangle=1\).  In the cylinder frame,
inserting a complete basis of states gives
\begin{align}
  2^{-4\Delta_\phi}\mathcal G(u,v)
  &=
  \bra{\vac}
    \widetilde\phi(0,0)\widetilde\phi(0,\pi)
    \widetilde\phi(\log r,\alpha)
    \widetilde\phi(\log r,\alpha+\pi)
  \ket{\vac}
  \notag\\
  &=
  \sum_n r^{E_n}
  \bra{\vac}
    \widetilde\phi(0,0)\widetilde\phi(0,\pi)
  \ket{n}
  \bra{n}
    \widetilde\phi(0,\alpha)\widetilde\phi(0,\alpha+\pi)
  \ket{\vac}.
  \label{eq:cylinder-four-point-complete-set}
\end{align}
For a primary \(\Phi_i\) of spin \(s_i\) and dimension \(\Delta_i\), its
level-\(n\) descendants have energy \(\Delta_i+n\).  If
\(\underline N=(\beta;N_1,\ldots,N_D)\), with
\(\sum_\mu N_\mu=n\), labels a descendant
\[
  \ket{\Phi_i,\underline N}
  =
  \prod_{\mu=1}^D \widehat P_\mu^{\,N_\mu}
  \ket{\Phi_i,\beta},
\]
then the Gram matrix
\[
  G_{\underline N\underline M}
  =
  \langle
    \Phi_i,\underline N\,|\,\Phi_i,\underline M
  \rangle
\]
and its inverse contract the two matrix elements in
\eqref{eq:cylinder-four-point-complete-set}.  If null descendants are
present, one first passes to a reduced linearly independent basis.  Therefore
a block has a radial expansion
\begin{equation}
  2^{-4\Delta_\phi}
  G_{\Delta_i,s_i}(u,v)
  =
  \sum_{n=0}^{\infty}
  A_{\Delta_i,s_i,n}(\alpha)\,r^{\Delta_i+n},
  \label{eq:block-radial-expansion}
\end{equation}
where each coefficient is fixed by conformal representation theory.

The leading coefficient is especially simple.  At level zero the only angular
dependence is the contraction of two symmetric traceless spin-\(s\) tensors
made from the two unit vectors separated by angle \(\alpha\).  Hence, with a
standard normalization of the block,
\begin{equation}
  G_{\Delta,s}(u,v)
  =
  r^\Delta
  C_s^{(D/2-1)}(\cos\alpha)
  +
  O(r^{\Delta+1}),
  \label{eq:block-leading-gegenbauer}
\end{equation}
where \(C_s^{(\lambda)}\) is the Gegenbauer polynomial.  For example this is
\(\cos(s\alpha)\) in \(D=2\), the Legendre polynomial \(P_s(\cos\alpha)\) in
\(D=3\), and \(\sin((s+1)\alpha)/((s+1)\sin\alpha)\) in \(D=4\), after the
corresponding conventional normalizations.

\section{Reflection Positivity and OPE Coefficients}

Reflection positivity constrains the numerical OPE data.  Choose Hermitian
scalar primaries with two-point normalization
\[
  \langle \mathcal O_A(x)\mathcal O_B(0)\rangle
  =
  \frac{\delta_{AB}}{|x|^{2\Delta_A}} .
\]
With this normalization, phases of nondegenerate Hermitian primaries can be
chosen so that three-point coefficients of Hermitian scalar operators are
real.  If \(\phi\) is a Hermitian scalar primary, the contribution of a
primary family \(\mathcal O\) to the \(\phi\phi\phi\phi\) four-point function
has coefficient
\[
  \lambda_{\phi\phi\mathcal O}^2
  \ge0
\]
in the channel where the first pair of operators is radially inside the
second pair.  This nonnegative number is the squared norm of the projection of
the state \(\phi(x)\phi(0)\vac\) onto the conformal multiplet of
\(\mathcal O\).

When several operators with the same quantum numbers mix, the statement is
matrix-valued.  Let \(P_\Delta\) be the spectral projection onto a finite
degenerate subspace of fixed spin and scaling dimension.  For any coefficients
\(\alpha_i\),
\[
  \left\|
    P_\Delta
    \sum_i\alpha_i\,\mathcal O_i(x_i)\vac
  \right\|^2
  \ge0 .
\]
The corresponding OPE coefficients form a positive semidefinite matrix in
that degenerate subspace.  Positivity of OPE data is therefore a direct
Hilbert-space consequence of radial quantization, not a separate dynamical
assumption.

\section{Associativity}

Consider a correlation function with three nearby operators.  There are two
ways to apply the OPE:
\[
  (\mathcal O_1\mathcal O_2)\mathcal O_3,
  \qquad
  \mathcal O_1(\mathcal O_2\mathcal O_3).
\]
Both compute the same local state on a sphere enclosing all three insertions.
Associativity of the OPE is the equality of these two expansions in their
common domain of convergence and by analytic continuation beyond it.  For
four-point functions this associativity becomes crossing symmetry.  In the
core monograph, crossing is used only at this structural level: a complete
solution theory for crossing equations requires additional analytic and
numerical machinery and belongs to a later dedicated treatment.

For identical scalar primaries \(\phi\) of dimension \(\Delta_\phi\), write
the four-point function in the standard form
\[
  \langle
    \phi(x_1)\phi(x_2)\phi(x_3)\phi(x_4)
  \rangle
  =
  \frac{\mathcal G(u,v)}
       {|x_{12}|^{2\Delta_\phi}|x_{34}|^{2\Delta_\phi}},
\]
where \(x_{ij}=x_i-x_j\) and the conformal cross-ratios are
\[
  u=\frac{x_{12}^2x_{34}^2}{x_{13}^2x_{24}^2},
  \qquad
  v=\frac{x_{14}^2x_{23}^2}{x_{13}^2x_{24}^2}.
\]
Exchanging \(x_2\) and \(x_4\) gives the crossing relation
\begin{equation}
  v^{\Delta_\phi}\mathcal G(u,v)
  =
  u^{\Delta_\phi}\mathcal G(v,u).
  \label{eq:identical-scalar-crossing}
\end{equation}
If \(\phi\) is Hermitian and the exchanged primaries are orthonormal, the
\(12\to34\) expansion takes the form
\begin{equation}
  \mathcal G(u,v)
  =
  1+
  \sum_{\mathcal O\ne\mathbf 1}
  \lambda_{\phi\phi\mathcal O}^2
  G_{\Delta_{\mathcal O},s_{\mathcal O}}(u,v),
  \qquad
  \lambda_{\phi\phi\mathcal O}^2\ge0.
  \label{eq:identical-scalar-positive-block-expansion}
\end{equation}
The identity block gives the constant \(1\).  The \(14\to23\) expansion gives
the same formula with \((u,v)\mapsto(v,u)\).  Equation
\eqref{eq:identical-scalar-crossing} is therefore not an additional
postulate; it is the equality of two convergent local expansions of one
Euclidean correlation function, continued to the same analytic function where
the continuations overlap.

\section{A First Consequence of Crossing and Positivity}

The simplest use of \eqref{eq:identical-scalar-crossing} is to see why a
unitary CFT cannot have an arbitrarily large gap above the identity in the
\(\phi\times\phi\) OPE while keeping \(\Delta_\phi\) fixed.  This subsection
is deliberately a controlled approximation, not a sharp bootstrap theorem:
we keep the leading radial behavior of heavy exchanged blocks near the
crossing-symmetric point and state the assumptions used.

Take \(z=\bar z\in(0,1)\), so that \(u=z^2\) and \(v=(1-z)^2\).  The crossing
equation becomes
\begin{equation}
  (1-z)^{2\Delta_\phi}\mathcal G(z)
  =
  z^{2\Delta_\phi}\mathcal G(1-z),
  \label{eq:real-line-identical-crossing}
\end{equation}
where \(\mathcal G(z)\) abbreviates \(\mathcal G(z^2,(1-z)^2)\).  Using
\eqref{eq:identical-scalar-positive-block-expansion},
\begin{equation}
  \bigl[(1-z)^{2\Delta_\phi}-z^{2\Delta_\phi}\bigr]
  +
  \sum_{\mathcal O\ne\mathbf1}
  \lambda_{\mathcal O}^2
  \left[
    (1-z)^{2\Delta_\phi}G_{\mathcal O}(z)
    -
    z^{2\Delta_\phi}G_{\mathcal O}(1-z)
  \right]
  =
  0.
  \label{eq:crossing-expanded-identity-plus-blocks}
\end{equation}
For \(0<z<1\), define
\[
  r(z)=\frac{z}{(1+\sqrt{1-z})^2}.
\]
If the exchanged dimensions contributing above the identity are large, and if
we keep only the leading radial behavior
\[
  G_{\mathcal O}(z)
  \simeq
  r(z)^{\Delta_{\mathcal O}}
\]
on this line, then \eqref{eq:crossing-expanded-identity-plus-blocks} is
approximated by
\begin{equation}
  \bigl[(1-z)^{2\Delta_\phi}-z^{2\Delta_\phi}\bigr]
  +
  \sum_{\mathcal O\ne\mathbf1}
  \lambda_{\mathcal O}^2
  \left[
    (1-z)^{2\Delta_\phi}r(z)^{\Delta_{\mathcal O}}
    -
    z^{2\Delta_\phi}r(1-z)^{\Delta_{\mathcal O}}
  \right]
  =
  0.
  \label{eq:large-gap-approx-crossing}
\end{equation}
The crossing-symmetric point is \(z=1/2\), where
\[
  r(1/2)=3-2\sqrt2.
\]
Set \(z=1/2+x\) and expand \eqref{eq:large-gap-approx-crossing} in \(x\).
The identity contribution is
\[
  (1-z)^{2\Delta_\phi}-z^{2\Delta_\phi}
  =
  -\Delta_\phi 2^{3-2\Delta_\phi}
  \left[
    x+\frac13(\Delta_\phi-1)(2\Delta_\phi-1)x^3+\cdots
  \right].
\]
For a heavy exchanged primary, the leading odd part of the approximate block
contribution has the form
\[
  \widetilde\lambda_{\mathcal O}^{\,2}
  \left[
    x+\frac43\Delta_{\mathcal O}^2x^3+\cdots
  \right],
  \qquad
  \widetilde\lambda_{\mathcal O}^{\,2}\ge0,
\]
where the positive factor relating
\(\widetilde\lambda_{\mathcal O}^{\,2}\) to
\(\lambda_{\mathcal O}^2\) includes
\((3-2\sqrt2)^{\Delta_{\mathcal O}}\) and other positive constants.  Matching
the coefficients of \(x\) and \(x^3\) gives, in this approximation,
\[
  \sum_{\mathcal O\ne\mathbf1}
  \widetilde\lambda_{\mathcal O}^{\,2}
  \left[
    \Delta_{\mathcal O}^2
    -
    (\Delta_\phi-1)(2\Delta_\phi-1)
  \right]
  =
  0.
\]
Since every coefficient in the sum is nonnegative, not all exchanged
dimensions can be much larger than
\[
  \sqrt{(\Delta_\phi-1)(2\Delta_\phi-1)}.
\]
The robust conclusion is qualitative: crossing plus reflection positivity
forces nontrivial operator content in the \(\phi\times\phi\) OPE at dimensions
controlled by \(\Delta_\phi\).  A rigorous numerical or analytic bound
requires the full conformal blocks and lies beyond this first approximation.
